% LaTeX file for Complainify Project Report
\documentclass[12pt,a4paper]{report}

% Packages
\usepackage[utf8]{inputenc}
\usepackage[T1]{fontenc}
\usepackage{times} % Times New Roman
\usepackage{geometry}
\geometry{margin=1in, top=1in, bottom=1in, left=1.5in, right=1in}
\usepackage{setspace}
\onehalfspacing % 1.5 line spacing
\usepackage{graphicx}
\usepackage{hyperref}
\usepackage{titlesec}
\usepackage{tocloft}
\usepackage{longtable}
\usepackage{array}
\usepackage{booktabs}
\usepackage{fancyhdr}
\usepackage{enumitem}
\usepackage{float}
\usepackage{caption}
\usepackage{indentfirst}
\usepackage{chngcntr}

\counterwithout{figure}{chapter}
\counterwithout{table}{chapter}

% Page style
\pagestyle{fancy}
\fancyhf{}
\fancyhead[R]{\thepage}
\renewcommand{\headrulewidth}{0pt}

% Chapter formatting
\titleformat{\chapter}[display]
  {\bfseries\centering\Large}
  {\chaptertitlename\ \thechapter}
  {0pt}
  {\MakeUppercase}

\titlespacing{\chapter}{0pt}{0pt}{20pt}

\titleformat{\section}
  {\bfseries\normalsize}
  {\thesection}
  {1em}
  {}

\titleformat{\subsection}
  {\bfseries\normalsize}
  {\thesubsection}
  {1em}
  {}

% Table of contents formatting
\renewcommand{\cfttoctitlefont}{\hfill\bfseries\Large}
\renewcommand{\cftaftertoctitle}{\hfill}
\renewcommand{\cftchapfont}{\bfseries}
\renewcommand{\cftchappagefont}{\bfseries}
\setlength{\cftbeforetoctitleskip}{0pt}

% List of Figures/Tables
\renewcommand{\cftloftitlefont}{\hfill\bfseries\Large}
\renewcommand{\cftafterloftitle}{\hfill}
\renewcommand{\cftlottitlefont}{\hfill\bfseries\Large}
\renewcommand{\cftafterlottitle}{\hfill}

\begin{document}

% ============================================================
% COVER PAGE
% ============================================================
\begin{titlepage}
\begin{center}

\vspace*{1cm}

{\Large \textbf{A Project Report\\[0.3cm] on\\[0.3cm]
Complainify}}

\vspace{1.5cm}

{\normalsize Submitted in partial fulfillment of the requirement of\\[0.2cm]
Project-VI (BIT356CO)\\[0.2cm]
Bachelor in Information Technology}

\vspace{1.5cm}

{\normalsize \textbf{Submitted To}\\[0.2cm]
Purbanchal University\\
Biratnagar, Nepal}

\vspace{0.8cm}

{\normalsize \textbf{Submitted By}\\[0.2cm]
Mridul Prakash Shah (350560)\\
Om Prakash Basnet (350562)\\
Samir Poudel (350574)}

\vspace{1.5cm}

\includegraphics[width=3cm]{kcc_logo.png}\\[0.5cm]
{\normalsize \textbf{KANTIPUR CITY COLLEGE}\\[0.1cm]
Putalisadak, Kathmandu\\[0.1cm]
October, 2026}

\vfill
\end{center}
\end{titlepage}

% ============================================================
% TITLE PAGE (with supervisor)
% ============================================================
\begin{titlepage}
\begin{center}

\vspace*{1cm}

{\Large \textbf{A Project Report\\[0.3cm] on\\[0.3cm]
Complainify}}

\vspace{1.5cm}

{\normalsize Submitted in partial fulfillment of the requirement of\\[0.2cm]
Project-VI (BIT356CO)\\[0.2cm]
Bachelor in Information Technology}

\vspace{1.5cm}

{\normalsize \textbf{Submitted To}\\[0.2cm]
Purbanchal University\\
Biratnagar, Nepal}

\vspace{0.8cm}

{\normalsize \textbf{Submitted By}\\[0.2cm]
Mridul Prakash Shah (350560)\\
Om Prakash Basnet (350562)\\
Samir Poudel (350574)}

\vspace{0.8cm}

{\normalsize \textbf{Project Supervisor}\\[0.2cm]
Er. Maryada Karki\\
Mr. Santosh Shah}

\vspace{1.5cm}

\includegraphics[width=3cm]{kcc_logo.png}\\[0.5cm]
{\normalsize \textbf{KANTIPUR CITY COLLEGE}\\[0.1cm]
Putalisadak, Kathmandu\\[0.1cm]
May, 2026}

\vfill
\end{center}
\end{titlepage}

% ============================================================
% ABSTRACT
% ============================================================
\chapter*{Abstract}
\addcontentsline{toc}{chapter}{Abstract}

Educational institutions often face difficulties in managing student complaints because traditional methods such as paper forms, emails, or manual records lead to delays, poor organization, and a lack of transparency in the complaint resolution process. To address these issues, this project proposes Student Complainify, an AI-Based Student Complaint Analysis and Decision Support System that provides a centralized and user-friendly platform for handling student grievances. Through this system, students can submit complaints related to academics, examinations, administration, infrastructure, or faculty issues at any time and from any location. The system uses Artificial Intelligence techniques, including sentiment analysis and automatic text classification, to analyze the complaint content, identify its category, and determine its urgency level. Based on the analysis, complaints are automatically assigned priorities such as High, Medium, or Low and forwarded to the appropriate authorities for faster action. The system also includes an administrative dashboard that provides real-time complaint statistics, tracking facilities, filtering options, and report generation features to support effective decision-making. Additionally, an integrated email notification service keeps students updated about the status of their complaints and alerts administrators when high-priority issues are reported. Developed using the Waterfall Software Development Life Cycle model and implemented with technologies such as HTML, CSS, JavaScript, Python Flask, MySQL, and machine learning libraries, the proposed system aims to improve efficiency, enhance transparency, and ensure quicker and more effective resolution of student complaints in educational institutions.

% ============================================================
% ACKNOWLEDGEMENT
% ============================================================
\chapter*{Acknowledgement}
\addcontentsline{toc}{chapter}{Acknowledgement}

We would like to express our heartfelt gratitude to everyone who has encouraged us to work on this project. First and foremost, we would like to express our gratitude to the entire team of Kantipur City College (KCC) for this opportunity, especially to the Principal of Kantipur City College, who assisted us in furthering our knowledge in this field.

We would like to express our special gratitude to our supervisors, Er. Maryada Karki and Mr. Santosh Shah, who have consistently encouraged, inspired, and provided us with a wealth of information that has been beneficial. Their advice was helpful in completing this assignment. We could not have asked for better supervisors, counselors, or mentors.

This project would not have been possible without the assistance of each team member, as well as the entire class, who offered suggestions, shared their experiences, and provided advice throughout the project. We are grateful for this.

Finally, we would like to express our gratitude to our friends and colleagues who supported us directly and indirectly throughout this project.

\vspace{2cm}
\begin{flushright}
Mridul Prakash Shah\\
Om Prakash Basnet\\
Samir Poudel
\end{flushright}

% ============================================================
% DECLARATION
% ============================================================
\chapter*{Declaration}
\addcontentsline{toc}{chapter}{Declaration}

We declare that this project report titled \textit{``Complainify''} submitted in partial fulfillment of the Bachelor of Information Technology is a record of original work carried out by us under the supervision of Er. Maryada Karki and Mr. Santosh Shah and has not formed the basis for the award of any other degree or Diploma, in this or any other institution or University. In keeping with the ethical practice of reporting scientific information, due acknowledgments have been made wherever the findings of others have been cited.

\vspace{2cm}
\begin{flushright}
Mridul Prakash Shah (350560)\\
Om Prakash Basnet (350562)\\
Samir Poudel (350574)
\end{flushright}

% ============================================================
% SUPERVISOR'S APPROVAL
% ============================================================
\chapter*{Supervisor's Approval}
\addcontentsline{toc}{chapter}{Supervisor's Approval}

This is to certify that the project entitled \textit{``Complainify''} undertaken and demonstrated by Mridul Prakash Shah, Om Prakash Basnet and Samir Poudel has been successfully completed under my supervision as a partial fulfillment of the requirements of the degree of BIT 6th semester under Purbanchal University. I, henceforth, approve this project to be awarded the certificate by the concerned authority. During supervision, I found students hardworking, skilled and ready to undertake any professional work related to this field in future.

\vspace{2cm}
\begin{minipage}{0.45\textwidth}
Mr. Santosh Shah\\
Project Supervisor\\
Department of IT\\
Date:
\end{minipage}
\hfill
\begin{minipage}{0.45\textwidth}
Er. Maryada Karki\\
Project Supervisor\\
Department of IT\\
Date:
\end{minipage}

% ============================================================
% CERTIFICATE FROM DEPARTMENT
% ============================================================
\chapter*{Certificate From Department}
\addcontentsline{toc}{chapter}{Certificate From Department}

Following the Supervisor's Approval and Examiner's Acceptance, the project entitled \textit{``Complainify''} submitted by Mridul Prakash Shah, Om Prakash Basnet and Samir Poudel as a partial fulfillment of the requirements for the degree of BIT 6th Semester under Purbanchal University, has been officially awarded by this certificate. I wish the students all the best for their future endeavors.

\vspace{2cm}
\begin{minipage}{0.45\textwidth}
Mr. Rubim Shrestha\\
Project Coordinator\\
Department of IT\\
Date:
\end{minipage}
\hfill
\begin{minipage}{0.45\textwidth}
Mr. Saroj Pandey\\
Head of Department\\
Department of IT\\
Date:
\end{minipage}

% ============================================================
% TABLE OF CONTENTS
% ============================================================
\renewcommand{\contentsname}{Table of Contents}
\tableofcontents
\newpage

% ============================================================
% LIST OF FIGURES
% ============================================================
\renewcommand{\listfigurename}{List of Figures}
\listoffigures
\newpage

% ============================================================
% LIST OF TABLES
% ============================================================
\renewcommand{\listtablename}{List of Tables}
\listoftables
\newpage

% ============================================================
% LIST OF ABBREVIATIONS
% ============================================================
\chapter*{List of Abbreviations}
\addcontentsline{toc}{chapter}{List of Abbreviations}

\begin{longtable}{p{3cm}p{10cm}}
API & Application Programming Interface \\
CPU & Central Processing Unit \\
CSS & Cascading Style Sheets \\
CRUD & Create, Read, Update, Delete \\
DB & Database \\
DFD & Data Flow Diagram \\
ER & Entity Relationship \\
Flask & Python Web Framework \\
HTML & HyperText Markup Language \\
HTTP & HyperText Transfer Protocol \\
IDE & Integrated Development Environment \\
IRR & Internal Rate of Return \\
JS & JavaScript \\
JSON & JavaScript Object Notation \\
ML & Machine Learning \\
MySQL & My Structured Query Language \\
Naive Bayes & Naive Bayes Classification Algorithm \\
NCF & Net Cash Flow \\
NLP & Natural Language Processing \\
NPR & Nepalese Rupee \\
NPV & Net Present Value \\
OTP & One Time Password \\
PBP & Payback Period \\
RAM & Random Access Memory \\
SDLC & Software Development Life Cycle \\
SMTP & Simple Mail Transfer Protocol \\
SQL & Structured Query Language \\
UI & User Interface \\
URL & Uniform Resource Locator \\
XAMPP & Cross-Platform Apache MySQL PHP Perl \\
\end{longtable}

% ============================================================
% CHAPTER 1: INTRODUCTION
% ============================================================
\chapter{Introduction}

\section{Overview}
Complainify is an AI-Based Student Complaint Analysis and Decision Support System developed to automate the management of student complaints in educational institutions. Traditional complaint handling methods, such as paper forms and manual records, often lead to delayed responses, improper categorization, and a lack of transparency. To overcome these issues, the proposed system provides a centralized web-based platform where students can submit and track complaints efficiently.

The system applies concepts from Artificial Intelligence and Data Mining, including Natural Language Processing (NLP), text classification, and sentiment analysis. When a complaint is submitted, the text is preprocessed and classified into categories such as Academic, Hostels, IT Support, Infrastructure, Financial Services, Administrative, Security, Maintenance, Transport, Canteen, Library, or Other using the Naive Bayes Classification Algorithm, which is implemented from scratch in Python due to its simplicity and effectiveness in text classification tasks.

The system also performs sentiment analysis using a lexicon-based approach to determine whether the complaint expresses a positive, neutral, or negative tone. Based on the analysis, complaints are assigned a confidence tier (Auto, Suggest, or Unknown) and priority is automatically boosted if negative sentiment is detected. Additionally, an email notification system is integrated to send complaint acknowledgements, status updates, and alerts for high-priority complaints to both students and administrators.

The system is implemented using HTML, CSS, JavaScript, Python Flask, and MySQL, along with supporting libraries such as Pandas, NumPy, and Matplotlib. The classification model is trained on a dataset of over 9,000 labeled complaints and achieves approximately 96\% accuracy on held-out test data. The proposed system aims to improve efficiency, transparency, and responsiveness in handling student grievances while demonstrating the practical application of Artificial Intelligence and Data Mining concepts studied in the sixth semester.

\section{Problem Statement}
Most educational institutions still manage student complaints manually through paper forms, emails, or unstructured communication channels. These traditional approaches suffer from several limitations:
\begin{itemize}[nolistsep]
    \item Complaints are not properly categorized.
    \item Urgent issues are often delayed due to the absence of prioritization mechanisms.
    \item Administrators lack analytical insights regarding complaint trends.
    \item Students do not receive timely updates regarding complaint status.
    \item There is no centralized platform to monitor and manage complaints efficiently.
\end{itemize}

Therefore, there is a need for an intelligent, automated, and centralized complaint management system that can analyze complaints, determine their priority, notify concerned authorities through email, and assist administrators in making data-driven decisions.

\section{Objectives}
The main objective of the project is:
\begin{itemize}[nolistsep]
    \item To develop a web-based AI-powered system for automated student complaint analysis, classification, sentiment detection, and priority assignment.
    \item To implement a Naive Bayes classifier from scratch for categorizing complaints into relevant departments.
    \item To integrate a lexicon-based sentiment analyzer for detecting the emotional tone of complaints.
    \item To provide role-based dashboards for students and administrators with real-time statistics and tracking.
    \item To generate analytical reports and visualizations for data-driven decision making.
\end{itemize}

\section{Features}
Complainify is designed with several features to ensure usability, efficiency, and reliability:

\begin{itemize}[nolistsep]
    \item \textbf{User Registration and Authentication:} Secure multi-role registration and login system for Students and Administrators, with email-based password reset and OTP verification features.
    
    \item \textbf{Complaint Submission System:} Students can submit complaints related to academics, hostels, IT, infrastructure, administration, or other issues by providing complaint details through an online portal.
    
    \item \textbf{AI-Based Complaint Classification:} The system automatically analyzes complaint text using a Multinomial Naive Bayes classifier implemented from scratch to classify complaints into appropriate categories with a confidence score.
    
    \item \textbf{Sentiment Analysis:} The system performs lexicon-based sentiment analysis on complaint content and identifies whether the complaint expresses positive, neutral, or negative sentiments to better understand the severity of issues.
    
    \item \textbf{Automatic Priority Assignment:} Based on complaint content, sentiment, and predefined rules, the system automatically assigns priority levels such as High, Medium, or Low. Negative sentiment complaints receive a priority boost.
    
    \item \textbf{Complaint Status Tracking:} Students can monitor the status of their complaints, including Pending, In Progress, or Resolved, ensuring transparency throughout the complaint resolution process.
    
    \item \textbf{Search and Filter Functionality:} Administrators can search and filter complaints based on category, priority, status, or date for efficient complaint management.
    
    \item \textbf{Dashboard and Analytics:} Role-specific dashboards provide complaint statistics, category-wise distributions, sentiment analysis, priority analysis, and graphical reports to support data-driven decision-making.
    
    \item \textbf{Email Notification System:} Automated email notifications are sent for complaint submission confirmations, status updates, and assignment notifications using SMTP services.
    
    \item \textbf{CSV Report Export:} Administrators can generate and export complaint reports in CSV format containing complaint summaries and statistics.
    
    \item \textbf{Three-Tier Confidence System:} Classifier outputs include a confidence tier — Auto (90\%+), Suggest (60-89\%), or Unknown (below 60\%) — helping administrators gauge classification reliability.
    
    \item \textbf{Responsive User Interface:} A clean, modern, and user-friendly interface is provided for both students and administrators, ensuring accessibility across desktop and mobile devices.
\end{itemize}

\section{Significance}
Complainify is significant because it provides an intelligent and centralized platform for managing student grievances efficiently by automating complaint submission, classification, prioritization, and tracking. The system improves transparency and communication between students and administrators through real-time status updates and email notifications while reducing manual effort and delays in complaint resolution.

By integrating Artificial Intelligence and Data Mining techniques such as text classification, sentiment analysis, and pattern discovery, the system helps institutions identify recurring issues, analyze complaint trends, and make informed decisions through analytical dashboards and reports. The project demonstrates the practical application of AI and Data Mining concepts in solving real-world problems and contributes to improving student satisfaction and institutional performance.

\section{Scope and Limitations}

\subsection{Scope}
The scope of Complainify is to develop a web-based intelligent complaint management system that allows students to submit and track complaints while enabling administrators to manage and analyze them efficiently. The system uses Artificial Intelligence and Data Mining techniques such as Natural Language Processing, Naive Bayes Classification, and sentiment analysis to categorize complaints, assign priorities, generate reports, and provide email notifications, thereby improving transparency and supporting data-driven decision-making in educational institutions.

\subsection{Limitations}
While developing the project, some features could not be implemented due to time and resource limitations:
\begin{itemize}[nolistsep]
    \item Integration with external university management systems is not implemented.
    \item Multi-language support is currently unavailable.
    \item Advanced AI-based recommendation and predictive complaint resolution features are not included.
    \item Real-time mobile push notification functionality is not implemented; the system currently relies on email notifications.
    \item Voice-based complaint submission and multilingual Natural Language Processing features are not included.
    \item The accuracy of complaint classification depends on the quality and quantity of training data available.
    \item The lexicon-based sentiment analyzer may not capture nuanced or sarcastic expressions accurately.
\end{itemize}

\section{Documentation Organization}
\textbf{Chapter 1} introduces our project by explaining its purpose, importance, objectives, features, scope, and limitations, providing an overview of the system. \textbf{Chapter 2} reviews existing complaint management systems, compares their strengths and limitations, and identifies gaps that justify the development of the proposed system. \textbf{Chapter 3} describes the methodology used to develop the project, including the Software Development Life Cycle, tools, technologies, programming languages, database systems, and team roles. \textbf{Chapter 4} focuses on system design and requirement analysis, including functional requirements, feasibility study, and system architecture. \textbf{Chapter 5} presents system diagrams such as Data Flow Diagrams, Entity Relationship Diagrams, Use Case Diagrams, and Gantt charts to illustrate system workflow and project planning. \textbf{Chapter 6} explains system development and implementation details, including hardware and software requirements. \textbf{Chapter 7} discusses system testing and debugging, presenting test cases and results. Finally, \textbf{Chapter 8} concludes the project with a summary of achievements and suggestions for future improvements.

% ============================================================
% CHAPTER 2: LITERATURE REVIEW
% ============================================================
\chapter{Literature Review}

Efficient complaint management is important in educational institutions to ensure that students' concerns are addressed effectively and within an appropriate timeframe. Traditional complaint handling approaches are mostly manual and often suffer from delayed responses, lack of categorization, and difficulties in identifying urgent issues. To overcome these problems, several intelligent complaint management and feedback analysis systems have been proposed. These systems employ machine learning, Natural Language Processing, sentiment analysis, and automated prioritization to improve complaint handling and support decision making.

Since the proposed system is developed as a web-based application with AI-assisted analysis, reviewing existing systems helps identify their features, strengths, and limitations. Most existing systems focus on complaint classification or sentiment analysis, while the proposed system combines sentiment analysis, automatic categorization, priority assignment, and dashboard-based reporting.

\section{Smart Complaint Management System with AI-Powered Prioritization and Escalation}
The Smart Complaint Management System with AI-Powered Prioritization and Escalation is an AI-based complaint handling system designed to improve the efficiency of complaint management through Natural Language Processing and automated prioritization. The system classifies complaints and determines their urgency to ensure timely responses and effective resolution \cite{deo2025}.

\subsection{Working Principle}
The system provides a web interface through which users can submit complaints. After submission, Natural Language Processing techniques are used to analyze the complaint text and classify it into appropriate categories. Sentiment analysis and AI-based priority assignment are used to determine the severity of complaints and automatically escalate urgent cases to the concerned authorities \cite{deo2025}.

\subsection{Advantages}
\begin{itemize}[nolistsep]
    \item Reduces manual effort in complaint categorization.
    \item Provides automatic priority assignment.
    \item Supports dashboard-based monitoring and reporting.
    \item Enhances transparency in complaint handling.
\end{itemize}

\subsection{Limitations}
\begin{itemize}[nolistsep]
    \item Performance depends on the quality of training data.
    \item Incorrect sentiment analysis may affect priority assignment.
    \item Requires periodic model updates.
    \item Limited support for multilingual complaints.
\end{itemize}

\section{Machine Learning-Based Smart Students' Complaint System}
The Machine Learning-Based Smart Students' Complaint System was developed to automate the handling of academic complaints and improve the efficiency of complaint resolution in educational institutions. The system employs machine learning and deep learning techniques to classify complaints according to departments and identify the corresponding issue categories \cite{javed2025}.

\subsection{Working Principle}
Students submit complaints through the system, and the complaint data are processed using Natural Language Processing techniques. Various algorithms including Support Vector Machine, Random Forest, Decision Tree, Convolutional Neural Network, and BERT are applied for department classification and aspect identification. The classified complaints are then forwarded to the appropriate department for further action \cite{javed2025}.

\subsection{Advantages}
\begin{itemize}[nolistsep]
    \item Automates complaint classification.
    \item Supports multiple machine learning algorithms.
    \item Improves efficiency of complaint resolution.
    \item Reduces dependency on manual processing.
    \item Provides department-wise categorization.
\end{itemize}

\subsection{Limitations}
\begin{itemize}[nolistsep]
    \item Accuracy depends on dataset quality.
    \item Limited ability to detect implicit meanings in complaints.
    \item Requires a large amount of training data.
\end{itemize}

\section{AI-Driven Real-Time Feedback System for Enhanced Student Support}
The AI-Driven Real-Time Feedback System for Enhanced Student Support is a machine learning-based framework that integrates sentiment analysis and predictive algorithms to provide personalized support and improve student engagement. The system analyzes student interactions and emotional states to generate context-sensitive feedback \cite{prakash2024}.

\subsection{Working Principle}
The system collects data from student interactions such as forum posts, assignment submissions, and feedback forms. Sentiment analysis is applied to identify emotional patterns, while machine learning algorithms such as Decision Trees, Support Vector Machines, Random Forest, and Deep Learning models are used to predict student engagement and performance. Based on the analysis, personalized recommendations and interventions are provided to support students academically and emotionally \cite{prakash2024}.

\subsection{Advantages}
\begin{itemize}[nolistsep]
    \item Provides real-time analysis of student feedback.
    \item Enhances student engagement and satisfaction.
    \item Supports personalized recommendations.
    \item Improves prediction accuracy using sentiment information.
    \item Enables data-driven decision making.
\end{itemize}

\subsection{Limitations}
\begin{itemize}[nolistsep]
    \item Deep learning models require high computational resources.
    \item May not generalize well across different institutional contexts.
\end{itemize}

\section{Summary of Literature Review}
The review of existing systems reveals that while several AI-based complaint management and feedback analysis systems have been developed, most focus either on classification or sentiment analysis independently. Few systems integrate both features along with automated priority assignment, role-based dashboards, and email notifications in a single platform. Complainify addresses this gap by combining Naive Bayes classification, lexicon-based sentiment analysis, a three-tier confidence system, and real-time analytics dashboards into one integrated solution. Additionally, the classifier is implemented entirely from scratch, demonstrating a deeper understanding of the underlying algorithms without reliance on external machine learning libraries.

% ============================================================
% CHAPTER 3: METHODOLOGY
% ============================================================
\chapter{Methodology}

\section{System Development Cycle}
For the development of Complainify, the Waterfall Model was selected as the software development methodology. The Waterfall Model is a structured and sequential approach that divides the development process into clearly defined phases. Each phase must be completed before moving to the next phase, ensuring systematic development, proper documentation, and effective project management.

The Waterfall Model was chosen for the following reasons:

\begin{itemize}[nolistsep]
    \item \textbf{Clear and Fixed Requirements:} The Complainify system has well-defined requirements such as user registration, complaint submission, complaint classification, sentiment analysis, priority assignment, dashboard generation, and email notifications. Since these requirements remain stable throughout development, the Waterfall Model is suitable for the project.
    
    \item \textbf{Easy to Manage:} The project follows a step-by-step development approach where each phase has specific tasks and deliverables. This makes it easier to manage, monitor progress, and ensure that all modules are completed systematically, which is ideal for a student-level project.
    
    \item \textbf{Better Documentation:} The Waterfall Model emphasizes documentation at every stage, including requirement specifications, system design diagrams, database design, implementation details, and testing reports. Proper documentation is essential for academic projects and future maintenance.
    
    \item \textbf{Suitable for Academic Projects:} The Waterfall Model is simple, easy to understand, and suitable for small development teams. It helps students complete the project within a limited academic time frame while maintaining a structured development process.
    
    \item \textbf{Predictable Timeline:} Since development follows a sequential process, it is easier to estimate the time required for each phase and complete the project according to the academic schedule.
\end{itemize}

\begin{figure}[H]
\centering
\framebox{\parbox{10cm}{\centering \textbf{Figure 3.1: Waterfall Model}\\[0.3cm]
Requirements $\rightarrow$ Design $\rightarrow$ Implementation $\rightarrow$ Testing $\rightarrow$ Deployment $\rightarrow$ Maintenance}}
\caption{Waterfall Model}
\end{figure}

\section{Tools and Technologies Used}
Complainify has been developed using a carefully selected set of web technologies and tools to ensure reliability, usability, and efficient performance. The system is designed as a web-based application that enables students to submit complaints and administrators to manage them effectively. The development utilizes core web technologies along with additional libraries and APIs to provide a complete complaint management solution.

\subsection{Programming Languages and Web Technologies}
\begin{itemize}[nolistsep]
    \item \textbf{HTML:} Used to structure the web pages of the Complainify system. It defines the layout of forms, dashboards, and complaint management interfaces, allowing users to easily enter and view complaint details in an organized manner.
    
    \item \textbf{CSS:} Used to design and style the user interface. It helps create a clean, readable, and user-friendly layout, ensuring that the system is visually simple and accessible for all users.
    
    \item \textbf{JavaScript:} Used to add interactivity and dynamic behavior to the system. It handles form validation, AJAX calls for AI category preview, and responsive features that improve user experience.
    
    \item \textbf{Python (Flask):} Used as the backend web framework. Flask handles routing, session management, database connectivity, API endpoints, and integration with the machine learning classifier and sentiment analyzer.
\end{itemize}

\subsection{Database Management System}
MySQL is used as the database management system for Complainify. It stores all system data such as user accounts, complaints with their categories and sentiments, assignment records, and OTP verification data. The database is designed with properly related tables to maintain data consistency and allow efficient storage and retrieval of complaint information.

\subsection{Additional Libraries and Tools}
Several additional tools were integrated to enhance system functionality:
\begin{itemize}[nolistsep]
    \item \textbf{Pandas and NumPy:} Used for data preprocessing, dataset loading, and numerical computations for the Naive Bayes classifier.
    \item \textbf{Matplotlib:} Used for generating visualizations and charts in the analysis notebooks.
    \item \textbf{pymysql:} Used as the MySQL database connector for Flask.
    \item \textbf{Jupyter Notebook:} Used for developing and documenting the Naive Bayes training pipeline and sentiment analysis.
    \item \textbf{Tailwind CSS:} Used for rapid UI development with utility-first CSS classes.
    \item \textbf{SMTP (smtplib):} Used for sending email notifications for complaint acknowledgements, status updates, and OTP verification.
\end{itemize}

\section{Assignment of Roles and Responsibilities}

\begin{table}[H]
\centering
\caption{Assignment of Roles and Responsibilities}
\begin{tabular}{p{4cm}p{10cm}}
\toprule
\textbf{Group Member} & \textbf{Roles and Responsibilities} \\
\midrule
\textbf{Mridul Prakash Shah}
(Backend Lead) &
\textbf{Coding:} Flask backend API — all routes (login, register, submit complaint, admin dashboard, etc.), MySQL database schema, authentication system, session management, Naive Bayes classifier implementation from scratch, sentiment analyzer integration, AI prediction API endpoint. \\
& \textbf{Documentation:} Methodology Chapter, Feasibility Study (Technical, Operational, Economic, Schedule), Database Design, ER Diagram description, Requirement Gathering table. \\
\midrule
\textbf{Om Prakash Basnet}
(Testing and Debugging) &
\textbf{Coding:} JavaScript frontend logic, AJAX integration for AI category preview, admin dashboard analytics, CSV export functionality, testing and debugging all modules end-to-end, classifier test cases and evaluation. \\
& \textbf{Documentation:} Literature Review, All Test Cases (Chapter 7), DFD Level 1, ER Diagram, Functional Requirements. \\
\midrule
\textbf{Samir Poudel}
(Frontend Lead) &
\textbf{Coding:} Landing page (index.html), Login and Registration UI, Forgot Password and Reset Password flows (OTP UI), Student Dashboard layout and sidebar navigation, Submit Complaint form, Admin Dashboard, Track Complaint page, Settings pages, All template designs. \\
& \textbf{Documentation:} System Architecture and n-Tier layer description, Use Case Diagram, Introduction — Overview, Problem Statement, Objectives, DFD Level 0 and Level 1, Final editing and formatting of project report. \\
\bottomrule
\end{tabular}
\end{table}

% ============================================================
% CHAPTER 4: SYSTEM DESIGN AND ANALYSIS
% ============================================================
\chapter{System Design and Analysis}

\section{Requirements Analysis}
Complainify was developed to address the growing need for efficient complaint management in educational institutions. The requirements were gathered through analysis of existing complaint handling methods, study of manual complaint management challenges faced by students and administrative staff, and review of existing digital complaint management systems.

\subsection{Requirement Gathering}
Requirements were gathered through analysis of student complaint handling challenges, identification of communication gaps between students and administration, and review of existing manual and digital complaint management practices. The gathered requirements emphasize role-based access control, automated complaint classification, sentiment analysis, real-time tracking, and email notifications. Table 4.1 summarizes the gathered requirements.

\begin{table}[H]
\centering
\caption{Requirement Gathering}
\begin{tabular}{cp{10cm}}
\toprule
\textbf{ID} & \textbf{Description} \\
\midrule
R1 & The system shall provide separate authentication mechanisms for Students and Administrators using email and password. \\
R2 & Students shall be able to register, login, submit complaints, and manage their own profile and account settings. \\
R3 & Students shall be able to submit complaints with subject, description, and optional category selection. \\
R4 & The system shall automatically classify complaints into appropriate categories using AI-based text classification. \\
R5 & The system shall perform sentiment analysis on complaint text to detect positive, neutral, or negative tone. \\
R6 & The system shall automatically assign priority (High, Medium, Low) based on complaint content and sentiment. \\
R7 & Students shall be able to track the status of their complaints using a unique ticket ID. \\
R8 & Administrators shall be able to view all complaints with filtering by status, category, and priority. \\
R9 & Administrators shall be able to assign complaints to departments and update complaint status. \\
R10 & Administrators shall be able to view analytics dashboards with complaint statistics and trends. \\
R11 & The system shall send email notifications for complaint submission, assignment, and status changes. \\
R12 & The system shall provide a forgot password mechanism with OTP verification via email. \\
R13 & The system shall support CSV export of complaint data for reporting purposes. \\
R14 & Administrators shall be able to validate complaints as legitimate. \\
R15 & The system shall maintain an audit trail of complaint status changes and administrative actions. \\
\bottomrule
\end{tabular}
\end{table}

\subsection{Functional Requirements}
The functional requirements define the specific behaviors and operations the system must perform to meet user needs and business objectives. Each functional requirement is mapped to one or more gathered requirements to ensure full coverage. Table 4.2 details all functional requirements for Complainify.

\begin{longtable}{cp{8cm}p{2cm}}
\caption{Functional Requirements}\\
\toprule
\textbf{ID} & \textbf{Description} & \textbf{Reference} \\
\midrule
\endfirsthead
\multicolumn{3}{c}{\tablename\ \thetable\ -- Continued from previous page} \\
\toprule
\textbf{ID} & \textbf{Description} & \textbf{Reference} \\
\midrule
\endhead
\bottomrule
\endfoot
F1 & User Authentication: The system shall validate login credentials for Students and Administrators using email and password with SHA-256 hashing. Role-based redirection ensures each user lands on the correct dashboard. & R1 \\
F2 & Student Registration: New students shall register with full name, email, phone, and password. The system validates email uniqueness before account creation. & R2 \\
F3 & Complaint Submission: Students shall submit complaints with subject, description, optional category selection, and priority selection. The system accepts the submission and generates a unique ticket ID. & R3 \\
F4 & AI Classification: The system shall automatically classify complaint text into categories (Academics, Hostels, IT Support, etc.) using a Multinomial Naive Bayes classifier with a three-tier confidence system. & R4 \\
F5 & Sentiment Analysis: The system shall analyze complaint text using a lexicon-based approach and return a sentiment label (Positive, Neutral, Negative) with a numeric score. & R5 \\
F6 & Priority Assignment: The system shall assign priority (Low, Medium, High) and apply an automatic boost if negative sentiment is detected. & R6 \\
F7 & Complaint Tracking: Students shall track their complaints using a ticket ID, viewing status, category, priority, assigned department, and admin notes. & R7 \\
F8 & Admin Complaint Management: Administrators shall view, filter by status/category, assign to departments, and update complaint status. & R8, R9 \\
F9 & Admin Dashboard: Administrators shall view real-time statistics including total complaints, resolved counts, average resolution time, critical complaints, category breakdown, and sentiment distribution. & R10 \\
F10 & Email Notifications: The system shall send email notifications for complaint submission acknowledgement, assignment, and status updates via SMTP. & R11 \\
F11 & Forgot Password with OTP: The system shall generate a 6-digit OTP, store it with a 10-minute expiry, send it via email, and verify it during password reset. & R12 \\
F12 & CSV Export: Administrators shall export complaint data as a CSV file containing all complaint fields. & R13 \\
F13 & Complaint Validation: Administrators shall mark complaints as validated after review. & R14 \\
F14 & Profile Management: Both students and administrators shall update their profile information and change passwords. & R2 \\
F15 & Password Change: Users shall change their password after validating the current password, with minimum length enforcement and confirmation matching. & R2 \\
\end{longtable}

\section{Feasibility Study}

\subsection{Technical Feasibility}
The technical feasibility analysis confirms that Complainify can be developed effectively with the technology stack and skills available to the development team. The system is built on widely adopted and well-documented technologies including HTML5, CSS3, JavaScript, Python Flask, and MySQL — all of which are free to use and supported by large developer communities. The team possesses the required knowledge of frontend development, backend API design, relational database management, and REST communication. The machine learning components (Naive Bayes classifier and sentiment analyzer) are implemented from scratch, demonstrating the team's understanding of the underlying algorithms. The system is designed to run on standard web hosting infrastructure without requiring specialized hardware. A modular architecture ensures that individual components such as the classifier engine, sentiment analyzer, notification system, and report generator can be developed, tested, and maintained independently. Therefore, Complainify is technically feasible.

\subsection{Operational Feasibility}
The operational feasibility analysis shows that Complainify will be readily accepted and effectively used by its two primary user groups — students and administrators. The system resolves critical pain points: students often face delays in complaint resolution, lack visibility into complaint status, and have no centralized platform to submit grievances. Administrators struggle with manual categorization, prioritization, and tracking of complaints. Complainify addresses each of these with intuitive role-specific dashboards, automated classification, priority assignment, and real-time tracking. The interface uses familiar web UI patterns, requiring minimal training for first-time users. The system is browser-based and does not require installation, making adoption straightforward across devices. Overall, Complainify fits naturally into existing institutional workflows and is operationally feasible.

\subsection{Economic Feasibility}
Complainify is economically feasible and demonstrates strong financial viability. The project requires an initial investment of NPR 150,000 covering labor, documentation, hardware, setup, training, and contingency. Annual operating costs range from NPR 25,000 in Year 1 to NPR 32,000 by Year 5, covering hosting, maintenance, and technical support. The system generates measurable annual benefits by reducing manual complaint processing time, improving administrative efficiency, and enhancing student satisfaction. Benefits grow from NPR 100,000 in Year 1 to NPR 160,000 by Year 5.

\begin{table}[H]
\centering
\caption{Initial Development Cost}
\begin{tabular}{lp{7cm}r}
\toprule
\textbf{S.No.} & \textbf{Cost Component} & \textbf{Estimated Cost (NPR)} \\
\midrule
1 & Labor Cost (Development) — 3 developers × 2 months × NPR 15,000 & 90,000 \\
2 & System Design and Documentation & 15,000 \\
3 & Hardware and Utilities (Laptop, hosting, internet, electricity) & 20,000 \\
4 & SMTP/Email Setup and Configuration & 5,000 \\
5 & Training and Deployment & 10,000 \\
6 & Contingency Cost & 10,000 \\
\midrule
\multicolumn{2}{r}{\textbf{Total Initial Investment}} & \textbf{150,000} \\
\bottomrule
\end{tabular}
\end{table}

\begin{table}[H]
\centering
\caption{Net Cash Flow Analysis}
\begin{tabular}{cccc}
\toprule
\textbf{Year} & \textbf{Annual Benefits (NPR)} & \textbf{Operating Costs (NPR)} & \textbf{Net Cash Flow (NPR)} \\
\midrule
1 & 100,000 & 25,000 & 75,000 \\
2 & 115,000 & 27,000 & 88,000 \\
3 & 130,000 & 29,000 & 101,000 \\
4 & 145,000 & 30,000 & 115,000 \\
5 & 160,000 & 32,000 & 128,000 \\
\bottomrule
\end{tabular}
\end{table}

\subsubsection{Payback Period (PBP)}
The Payback Period determines how quickly the initial investment is recovered.

\textbf{Formula:} PBP = Year before recovery + (Remaining Investment / Next Year NCF)

Cumulative Cash Flow:
\begin{itemize}[nolistsep]
    \item Year 1 = 75,000
    \item Year 2 = 163,000
\end{itemize}

\textbf{Calculation:}
PBP = 1 + (150,000 - 75,000) / 88,000 = 1.85 years ≈ 1 year 10 months

The investment is recovered in less than 2 years, indicating low financial risk.

\subsubsection{Net Present Value (NPV)}
\textbf{Formula:} PV = NCF / (1+r)$^t$, NPV = $\sum$PV - Initial Investment (Discount Rate = 12\%)

\begin{table}[H]
\centering
\caption{Calculation of NPV}
\begin{tabular}{cccc}
\toprule
\textbf{Year} & \textbf{NCF (NPR)} & \textbf{Present Value (NPR)} \\
\midrule
1 & 75,000 & 66,964 \\
2 & 88,000 & 70,153 \\
3 & 101,000 & 71,887 \\
4 & 115,000 & 73,084 \\
5 & 128,000 & 72,642 \\
\midrule
\multicolumn{2}{r}{\textbf{NPV = 66,964 + 70,153 + 71,887 + 73,084 + 72,642 - 150,000}} & \textbf{204,730} \\
\bottomrule
\end{tabular}
\end{table}

Positive NPV of NPR 204,730 confirms the project is financially viable.

\subsubsection{Internal Rate of Return (IRR)}
The IRR is the discount rate at which NPV = 0. Using financial calculation methods:
IRR ≈ 52.4\%

This is significantly higher than the required return (12-15\%), indicating high profitability.

\begin{table}[H]
\centering
\caption{Summary of Financial Metrics}
\begin{tabular}{lr}
\toprule
\textbf{Metric} & \textbf{Value} \\
\midrule
Payback Period & 1 year 10 months \\
Net Present Value & NPR 204,730 \\
Internal Rate of Return & 52.4\% \\
\bottomrule
\end{tabular}
\end{table}

The economic analysis shows that Complainify is financially sound with a short payback period, high NPV, and strong IRR, making it a highly feasible and sustainable solution.

\subsection{Schedule Feasibility}
The schedule feasibility analysis confirms that the project can be delivered within a planned 16-week timeframe, spanning from February through May. The project begins with discussion and initiation (February W1-W3), followed by requirement analysis (February W2-W4) and system design (February W3-March W1). Frontend and backend development are carried out in parallel from March W2 to May W1, while testing runs from March W4 to May W1. Bug fixing is performed from April W2 to May W4, and documentation is maintained continuously from February W1 to May W4. The overlapping activities and proper task distribution confirm that the project is schedule feasible within the proposed timeframe.

\begin{figure}[H]
\centering
\framebox{\parbox{12cm}{\centering \textbf{Gantt Chart}\\[0.3cm]
Discussion: Feb W1-W3 \\
Requirement Analysis: Feb W2-W4 \\
System Design: Feb W3-Mar W1 \\
Frontend Development: Mar W2-May W1 \\
Backend Development: Mar W2-May W1 \\
Testing: Mar W4-May W1 \\
Bug Fixing: Apr W2-May W4 \\
Documentation: Feb W1-May W4}}
\caption{Gantt Chart}
\end{figure}

% ============================================================
% CHAPTER 5: SYSTEM DESIGN
% ============================================================
\chapter{System Design}

\section{System Architecture}
Complainify follows a multi-tier architecture that separates the presentation layer, business logic layer, data access layer, and external integration layer. This architectural pattern ensures modularity, maintainability, and scalability of the system. Each layer has a well-defined responsibility and communicates with adjacent layers through clearly defined interfaces, preventing tight coupling and enabling independent modification of any layer.

\begin{figure}[H]
\centering
\framebox{\parbox{12cm}{\centering \textbf{Figure 5.1: System Architecture}\\[0.3cm]
\textbf{Presentation Layer:} HTML, CSS, JavaScript, Tailwind CSS (Browser)\\[0.2cm]
$\downarrow$ HTTP Requests / $\uparrow$ HTML Responses\\[0.2cm]
\textbf{Business Logic Layer:} Python Flask (app.py)\\[0.2cm]
$\downarrow$ SQL Queries / $\uparrow$ Result Sets\\[0.2cm]
\textbf{Data Access Layer:} MySQL + pymysql\\[0.2cm]
$\downarrow$ API Calls / $\uparrow$ Responses\\[0.2cm]
\textbf{External Integration Layer:} SMTP (Email)}}
\caption{System Architecture}
\end{figure}

\subsection{Architectural Components}

\begin{table}[H]
\centering
\caption{System Architecture Components}
\begin{tabular}{p{3.5cm}p{11cm}}
\toprule
\textbf{Layer} & \textbf{Description} \\
\midrule
\textbf{Presentation Layer} (HTML, CSS, JavaScript, Tailwind CSS) & Provides role-based interfaces for Students and Administrators. Handles user input validation, form interactions, and AJAX calls for AI category preview. Includes all dashboard pages, auth screens, complaint forms, and settings. \\
\midrule
\textbf{Business Logic Layer} (Python Flask) & Contains core application logic, role enforcement, complaint classification using Naive Bayes, sentiment analysis using lexicon-based approach, priority computation, and notification triggering. Processes user requests and performs data validation before calling the database layer. \\
\midrule
\textbf{Data Access Layer} (MySQL + pymysql) & Handles all database interactions. Manages CRUD operations for users, complaints, and OTP records. Ensures data integrity through foreign key constraints and proper indexing. \\
\midrule
\textbf{External Integration Layer} (SMTP) & Delivers automated email notifications for complaint acknowledgements, status updates, assignment notifications, and OTP verification. Decoupled from core logic so notification channels can be modified independently. \\
\bottomrule
\end{tabular}
\end{table}

\section{Procedure Oriented Diagram}

\subsection{Context Diagram (DFD Level 0)}

\begin{figure}[H]
\centering
\framebox{\parbox{12cm}{\centering \textbf{DFD Level 0: Context Diagram}\\[0.3cm]
\begin{tabular}{c}
\textbf{Student} \\ $\longleftrightarrow$ \\ Complainify System \\ $\longleftrightarrow$ \\ \textbf{Administrator}
\end{tabular}\\[0.3cm]
Student: Submit Complaint, Track Status, Register, Login, Update Profile\\
Administrator: Manage Complaints, View Dashboard, Assign, Resolve, Export Reports}}
\caption{DFD Level 0 — Context Diagram}
\end{figure}

\subsection{Data Flow Diagram (Level 1)}

\begin{figure}[H]
\centering
\framebox{\parbox{12cm}{\centering \textbf{DFD Level 1}\\[0.3cm]
\begin{tabular}{ll}
\textbf{Processes:} & \\
1.0 & User Authentication (Login/Register/Forgot Password) \\
2.0 & Complaint Submission with AI Classification \\
3.0 & Sentiment Analysis and Priority Assignment \\
4.0 & Complaint Tracking \\
5.0 & Admin Complaint Management \\
6.0 & Dashboard and Analytics Generation \\
7.0 & Email Notification \\
8.0 & Report Export \\[0.3cm]
\textbf{Data Stores:} & Users, Complaints, OTPs \\[0.3cm]
\textbf{External Entities:} Student, Administrator
\end{tabular}}
\caption{DFD Level 1}
\end{figure}

\subsection{Use Case Diagram}

\begin{figure}[H]
\centering
\framebox{\parbox{12cm}{\centering \textbf{Use Case Diagram}\\[0.3cm]
\textbf{Student:} Register, Login, Submit Complaint, Track Complaint, View Dashboard, Update Profile, Forgot Password, Change Password\\[0.3cm]
\textbf{Administrator:} Login, View Dashboard, View All Complaints, Filter Complaints, Assign Complaint, Update Status, Validate Complaint, Resend Email, Export CSV, View Reports, Manage Profile, Change Password}}
\caption{Use Case Diagram}
\end{figure}

\section{Database Design}

\subsection{E-R Diagram}
The database for Complainify consists of three main tables: \texttt{users}, \texttt{complaints}, and \texttt{otps}. The \texttt{users} table stores student and administrator account information including fullname, email, phone, password (SHA-256 hashed), plain password (for reference), and role. The \texttt{complaints} table stores all complaint data including ticket ID, category, priority, status, subject, description, sentiment, sentiment score, assigned department, assignment timestamp, resolution timestamp, admin notes, and validation status. The \texttt{otps} table stores OTP records for the forgot-password flow with email, OTP code, expiry timestamp, and usage status.

\begin{figure}[H]
\centering
\framebox{\parbox{12cm}{\centering \textbf{ER Diagram}\\[0.3cm]
\texttt{users (id, fullname, email, phone, plain\_password, password, role, created\_at)}\\[0.2cm]
1 --- has many ---> * \texttt{complaints (user\_id FK)}\\[0.2cm]
\texttt{complaints (id, ticket\_id, user\_id FK, fullname, email, category, priority, subject, description, status, sentiment, sentiment\_score, assigned\_to, assigned\_at, resolved\_at, admin\_notes, validated, email\_sent, created\_at)}\\[0.2cm]
\texttt{otps (id, email, otp, created\_at, expires\_at, used)}}
\caption{ER Diagram}
\end{figure}

\section{Role-Based Design}

\begin{table}[H]
\centering
\caption{Role-Based Access Control Matrix}
\begin{tabular}{p{6cm}ccc}
\toprule
\textbf{Feature / Action} & \textbf{Student} & \textbf{Admin} \\
\midrule
Register / Login & Yes & Yes \\
Submit Complaint & Yes & No \\
View Own Complaints & Yes & No \\
Track Complaint by Ticket ID & Yes & Yes \\
View All Complaints & No & Yes \\
Filter/Search Complaints & No & Yes \\
Assign Complaint to Department & No & Yes \\
Update Complaint Status & No & Yes \\
Validate Complaint & No & Yes \\
View Analytics Dashboard & No & Yes \\
View Reports & No & Yes \\
Export CSV & No & Yes \\
Resend Email Notification & No & Yes \\
Update Profile & Yes & Yes \\
Change Password & Yes & Yes \\
Forgot Password / Reset Password & Yes & Yes \\
Receive Email Notifications & Yes & Yes \\
\bottomrule
\end{tabular}
\end{table}

% ============================================================
% CHAPTER 6: SYSTEM DEVELOPMENT AND IMPLEMENTATION
% ============================================================
\chapter{System Development and Implementation}

\section{Programming Platforms}

\begin{enumerate}[nolistsep]
    \item \textbf{Frontend Technologies: HTML, CSS, JavaScript, Tailwind CSS}
    
    Core technologies used to design and develop the user interface. HTML structures the web pages, CSS styles the layout, JavaScript adds interactivity and AJAX functionality, and Tailwind CSS provides a utility-first framework for rapid UI development. Together they create a modern and user-friendly web interface.
    
    \item \textbf{Backend Framework: Python Flask}
    
    Used to handle server-side logic and system functionality. Flask manages user authentication, session management, complaint submission, AI classification, sentiment analysis, database operations, and API endpoints. It provides secure communication between the frontend and database with Jinja2 templating for dynamic HTML rendering.
    
    \item \textbf{Machine Learning: Naive Bayes (From Scratch)}
    
    The Multinomial Naive Bayes classifier is implemented entirely from scratch without using sklearn or any ML library. It includes a custom stemmer with 25 rules, bigram feature extraction, Laplace smoothing (alpha=1.0), and a three-tier confidence system. The model is trained on 7,204 labeled complaints and achieves approximately 96\% accuracy on 1,806 test samples.
    
    \item \textbf{Sentiment Analysis: Lexicon-Based (From Scratch)}
    
    The sentiment analyzer uses a manually curated lexicon of approximately 150 words with intensity scores. It includes negation handling (words like "not" flip polarity) and intensifier amplification (words like "very" boost intensity). No external libraries like NLTK or VADER are used.
    
    \item \textbf{Development Environment: Visual Studio Code}
    
    Used for writing, editing, and testing the project code. Provides useful extensions, debugging tools, and integrated terminal support for Git version control.
    
    \item \textbf{Local Server Environment: XAMPP}
    
    Used to run Apache server and MySQL database locally. Allows testing of the web application in a local environment before deployment.
    
    \item \textbf{Database System: MySQL}
    
    Used to store user data, complaint records, and OTP verification data. Ensures reliable data storage and retrieval with proper indexing on ticket IDs and email fields.
    
    \item \textbf{Additional Libraries and Tools}
    
    \begin{itemize}[nolistsep]
        \item \textbf{Pandas/NumPy:} Data loading and numerical computations for ML training.
        \item \textbf{Matplotlib:} Chart and visualization generation in analysis notebooks.
        \item \textbf{pymysql:} MySQL database connector for Flask.
        \item \textbf{Jupyter Notebook:} ML training pipeline and sentiment analysis documentation.
        \item \textbf{smtplib:} SMTP-based email notifications.
    \end{itemize}
\end{enumerate}

\section{Hardware and Software Requirements}

\subsection{Hardware Requirements}

\textbf{Minimum Requirements:}
\begin{itemize}[nolistsep]
    \item Processor: Dual-Core CPU (Intel i3 or equivalent)
    \item RAM: 4 GB
    \item Storage: 500 MB free space
    \item Display: 1366 $\times$ 768 resolution
\end{itemize}

\textbf{Recommended Requirements:}
\begin{itemize}[nolistsep]
    \item Processor: Intel Core i5 or higher
    \item RAM: 8 GB
    \item Storage: 1 GB free space
    \item Display: Full HD (1920 $\times$ 1080)
\end{itemize}

\subsection{Software Requirements}
\begin{itemize}[nolistsep]
    \item Operating System: Windows / Linux / macOS
    \item Web Browser: Google Chrome or any modern browser
    \item Code Editor: Visual Studio Code
    \item Server Environment: XAMPP (Apache + MySQL)
    \item Database: MySQL 8.0+
    \item Backend Runtime: Python 3.10+
    \item Python Packages: Flask, pymysql, pandas, numpy, matplotlib, jupyter
    \item Frontend: HTML5, CSS3, JavaScript (ES6+), Tailwind CSS
\end{itemize}

% ============================================================
% CHAPTER 7: TESTING AND DEBUGGING
% ============================================================
\chapter{Testing and Debugging}

\section{Testing Approach}
Complainify was tested using a multi-level testing strategy covering unit testing, integration testing, and user acceptance testing. Unit tests verified the correctness of individual components including the Naive Bayes classifier, sentiment analyzer, and API endpoints in isolation. Integration tests confirmed that frontend pages correctly communicate with backend endpoints and that data flows as expected from user input through to database storage. User acceptance testing was performed by simulating real-world usage scenarios for each of the two user roles — Student and Administrator — to confirm that the system behaves correctly from an end-user perspective.

\section{Testing Tools}

\begin{table}[H]
\centering
\caption{Testing Tools}
\begin{tabular}{c p{10cm}}
\toprule
\textbf{S.N.} & \textbf{Tool / Specification} \\
\midrule
1 & Hardware: Laptop with 8 GB RAM, Intel Core i5 processor \\
2 & Web Browser: Google Chrome and Mozilla Firefox \\
3 & XAMPP / Apache: Local web server environment \\
4 & MySQL 8.0: Database verification via phpMyAdmin \\
5 & Python: Direct testing of classifier and sentiment modules \\
6 & Jupyter Notebook: Evaluation with confusion matrix, precision, recall, F1 \\
\bottomrule
\end{tabular}
\end{table}

\section{Test Cases}
The following tables document test cases organized by module. Each test case records the test ID, description, input conditions, expected output, actual output observed during testing, and the final status (Pass or Fail).

\subsection{Login Module}

\begin{longtable}{c p{3.5cm} p{3cm} p{3.5cm} p{3cm} c}
\caption{Login Module Test Cases}\\
\toprule
\textbf{TC ID} & \textbf{Description} & \textbf{Input} & \textbf{Expected Output} & \textbf{Actual Output} & \textbf{Status} \\
\midrule
\endfirsthead
\multicolumn{6}{c}{\tablename\ \thetable\ -- Continued} \\
\toprule
\textbf{TC ID} & \textbf{Description} & \textbf{Input} & \textbf{Expected Output} & \textbf{Actual Output} & \textbf{Status} \\
\midrule
\endhead
\bottomrule
\endfoot
1.1 & Valid student login & Correct email and password (student) & Student dashboard opens & Student dashboard opens & Pass \\
1.2 & Valid admin login & Correct email and password (admin) & Admin dashboard opens & Admin dashboard opens & Pass \\
1.3 & Invalid email & Unregistered email & Error: Account not found & Error displayed & Pass \\
1.4 & Wrong password & Valid email, wrong password & Error: Incorrect password & Error displayed & Pass \\
1.5 & Empty email field & Blank email & Validation error & Error displayed & Pass \\
1.6 & Empty password field & Blank password & Validation error & Error displayed & Pass \\
1.7 & Password less than 6 chars & 3-character password & Error: Password too short & Error displayed & Pass \\
\end{longtable}

\subsection{Sign Up Module}

\begin{longtable}{c p{3.5cm} p{3cm} p{3.5cm} p{3cm} c}
\caption{Sign Up Module Test Cases}\\
\toprule
\textbf{TC ID} & \textbf{Description} & \textbf{Input} & \textbf{Expected Output} & \textbf{Actual Output} & \textbf{Status} \\
\midrule
\endfirsthead
\multicolumn{6}{c}{\tablename\ \thetable\ -- Continued} \\
\toprule
\textbf{TC ID} & \textbf{Description} & \textbf{Input} & \textbf{Expected Output} & \textbf{Actual Output} & \textbf{Status} \\
\midrule
\endhead
\bottomrule
\endfoot
2.1 & Valid student registration & All fields valid & Account created, redirected to login & Account created & Pass \\
2.2 & Duplicate email & Already registered email & Error: Email already registered & Error displayed & Pass \\
2.3 & Missing full name & Blank name & Validation error & Error displayed & Pass \\
2.4 & Invalid email format & Email without @ & Error: Valid email required & Error displayed & Pass \\
2.5 & Password mismatch & Different passwords & Error: Passwords do not match & Error displayed & Pass \\
2.6 & Password too short & Less than 6 chars & Error: Minimum 6 characters & Error displayed & Pass \\
2.7 & Admin registration & Valid admin details & Admin account created & Account created & Pass \\
\end{longtable}

\subsection{Forgot Password / Reset Password Module}

\begin{longtable}{c p{3.5cm} p{3cm} p{3.5cm} p{3cm} c}
\caption{Forgot Password and Reset Password Module Test Cases}\\
\toprule
\textbf{TC ID} & \textbf{Description} & \textbf{Input} & \textbf{Expected Output} & \textbf{Actual Output} & \textbf{Status} \\
\midrule
\endfirsthead
\multicolumn{6}{c}{\tablename\ \thetable\ -- Continued} \\
\toprule
\textbf{TC ID} & \textbf{Description} & \textbf{Input} & \textbf{Expected Output} & \textbf{Actual Output} & \textbf{Status} \\
\midrule
\endhead
\bottomrule
\endfoot
3.1 & Valid email OTP sent & Registered email & OTP stored in DB, redirect to reset & OTP sent, redirected & Pass \\
3.2 & Unregistered email & Email not in system & Error: No account found & Error displayed & Pass \\
3.3 & Empty email field & Blank email & Validation error & Error displayed & Pass \\
3.4 & Valid OTP + new password & Correct OTP, new password & Password updated, redirected to login & Password updated & Pass \\
3.5 & Wrong OTP entered & Incorrect 6-digit OTP & Error: Invalid or expired OTP & Error displayed & Pass \\
3.6 & New passwords don't match & Different new passwords & Error: Passwords do not match & Error displayed & Pass \\
\end{longtable}

\subsection{Complaint Submission Module}

\begin{longtable}{c p{3.5cm} p{3cm} p{3.5cm} p{3cm} c}
\caption{Complaint Submission Module Test Cases}\\
\toprule
\textbf{TC ID} & \textbf{Description} & \textbf{Input} & \textbf{Expected Output} & \textbf{Actual Output} & \textbf{Status} \\
\midrule
\endfirsthead
\multicolumn{6}{c}{\tablename\ \thetable\ -- Continued} \\
\toprule
\textbf{TC ID} & \textbf{Description} & \textbf{Input} & \textbf{Expected Output} & \textbf{Actual Output} & \textbf{Status} \\
\midrule
\endhead
\bottomrule
\endfoot
4.1 & Submit complaint with auto category & Subject + description, category=auto & Complaint saved with AI category & Category auto-detected & Pass \\
4.2 & Submit complaint with manual category & Select specific category & Complaint saved with selected category & Category applied & Pass \\
4.3 & Missing subject & Blank subject field & Validation error & Error displayed & Pass \\
4.4 & Missing description & Blank description & Validation error & Error displayed & Pass \\
4.5 & Sentiment detected as Negative & Angry/frustrated text & Sentiment=Negative, priority boosted & Negative detected, priority up & Pass \\
4.6 & Sentiment detected as Positive & Appreciative text & Sentiment=Positive & Positive detected & Pass \\
4.7 & AI preview via AJAX & Type in description field & Preview shows AI-detected category & Preview shown & Pass \\
4.8 & High priority submission & Select priority=High & Complaint saved with High priority & Priority saved & Pass \\
\end{longtable}

\subsection{Complaint Tracking Module}

\begin{longtable}{c p{3.5cm} p{3cm} p{3.5cm} p{3cm} c}
\caption{Complaint Tracking Module Test Cases}\\
\toprule
\textbf{TC ID} & \textbf{Description} & \textbf{Input} & \textbf{Expected Output} & \textbf{Actual Output} & \textbf{Status} \\
\midrule
\endfirsthead
\multicolumn{6}{c}{\tablename\ \thetable\ -- Continued} \\
\toprule
\textbf{TC ID} & \textbf{Description} & \textbf{Input} & \textbf{Expected Output} & \textbf{Actual Output} & \textbf{Status} \\
\midrule
\endhead
\bottomrule
\endfoot
5.1 & Track valid ticket ID & Enter existing ticket ID & Complaint details displayed & Details shown & Pass \\
5.2 & Track invalid ticket ID & Enter non-existent ID & No results, empty state shown & No results & Pass \\
5.3 & Track with empty field & Submit blank & Validation error & Error displayed & Pass \\
\end{longtable}

\subsection{Admin Complaint Management Module}

\begin{longtable}{c p{3cm} p{3cm} p{3.5cm} p{3cm} c}
\caption{Admin Complaint Management Module Test Cases}\\
\toprule
\textbf{TC ID} & \textbf{Description} & \textbf{Input} & \textbf{Expected Output} & \textbf{Actual Output} & \textbf{Status} \\
\midrule
\endfirsthead
\multicolumn{6}{c}{\tablename\ \thetable\ -- Continued} \\
\toprule
\textbf{TC ID} & \textbf{Description} & \textbf{Input} & \textbf{Expected Output} & \textbf{Actual Output} & \textbf{Status} \\
\midrule
\endhead
\bottomrule
\endfoot
6.1 & View all complaints & Login as admin, navigate to complaints & All complaints listed & All complaints shown & Pass \\
6.2 & Filter by status & Select "Pending" filter & Only pending complaints shown & Filtered correctly & Pass \\
6.3 & Filter by category & Select category filter & Only matching category shown & Filtered correctly & Pass \\
6.4 & Assign complaint & Select department, click Assign & Status = In Progress, assigned & Assigned successfully & Pass \\
6.5 & Update status to In Progress & Select status, add notes & Status updated, notes saved & Updated & Pass \\
6.6 & Update status to Resolved & Select Resolved, add notes & Resolved with timestamp & Resolved & Pass \\
6.7 & Validate complaint & Click Validate button & validated = 1 & Validated & Pass \\
6.8 & View complaint detail & Click on ticket ID & Full complaint details shown & Details shown & Pass \\
\end{longtable}

\subsection{Admin Dashboard Module}

\begin{longtable}{c p{3.5cm} p{3cm} p{3.5cm} p{3cm} c}
\caption{Admin Dashboard Module Test Cases}\\
\toprule
\textbf{TC ID} & \textbf{Description} & \textbf{Input} & \textbf{Expected Output} & \textbf{Actual Output} & \textbf{Status} \\
\midrule
\endfirsthead
\multicolumn{6}{c}{\tablename\ \thetable\ -- Continued} \\
\toprule
\textbf{TC ID} & \textbf{Description} & \textbf{Input} & \textbf{Expected Output} & \textbf{Actual Output} & \textbf{Status} \\
\midrule
\endhead
\bottomrule
\endfoot
7.1 & View dashboard stats & Login as admin & Total, resolved, pending, critical counts shown & Live data displayed & Pass \\
7.2 & View category breakdown & Check category section & Category percentages from real DB data & Live breakdown shown & Pass \\
7.3 & View sentiment distribution & Check sentiment section & Positive/Neutral/Negative counts shown & Live counts shown & Pass \\
7.4 & View average resolution time & Check metric card & Real avg from DB calculation & Live avg shown & Pass \\
\end{longtable}

\subsection{Email Notification Module}

\begin{longtable}{c p{3.5cm} p{3cm} p{3.5cm} p{3cm} c}
\caption{Email Notification Module Test Cases}\\
\toprule
\textbf{TC ID} & \textbf{Description} & \textbf{Input} & \textbf{Expected Output} & \textbf{Actual Output} & \textbf{Status} \\
\midrule
\endfirsthead
\multicolumn{6}{c}{\tablename\ \thetable\ -- Continued} \\
\toprule
\textbf{TC ID} & \textbf{Description} & \textbf{Input} & \textbf{Expected Output} & \textbf{Actual Output} & \textbf{Status} \\
\midrule
\endhead
\bottomrule
\endfoot
8.1 & Email on complaint submission & Submit complaint with valid email & Acknowledgement email sent & Email sent (or logged if no SMTP) & Pass \\
8.2 & Email on status change & Admin updates status & Status update email sent & Email sent & Pass \\
8.3 & Email on assignment & Admin assigns complaint & Assignment notification sent & Email sent & Pass \\
8.4 & Email with no SMTP config & No SMTP env vars set & Email skipped with console log & Graceful skip & Pass \\
8.5 & Resend email & Admin clicks Resend & Email resent & Email sent & Pass \\
\end{longtable}

\subsection{Classifier and Sentiment Module}

\begin{longtable}{c p{3.5cm} p{3cm} p{3.5cm} p{3cm} c}
\caption{Classifier and Sentiment Module Test Cases}\\
\toprule
\textbf{TC ID} & \textbf{Description} & \textbf{Input} & \textbf{Expected Output} & \textbf{Actual Output} & \textbf{Status} \\
\midrule
\endfirsthead
\multicolumn{6}{c}{\tablename\ \thetable\ -- Continued} \\
\toprule
\textbf{TC ID} & \textbf{Description} & \textbf{Input} & \textbf{Expected Output} & \textbf{Actual Output} & \textbf{Status} \\
\midrule
\endhead
\bottomrule
\endfoot
9.1 & Naive Bayes accuracy & Run classifier on test set (1806 samples) & Accuracy \textgreater 90\% & 96\% accuracy & Pass \\
9.2 & Auto confidence tier & Complaint with clear category keywords & tier=auto, confidence \textgreater= 0.90 & auto tier & Pass \\
9.3 & Suggest confidence tier & Ambiguous complaint text & tier=suggest, confidence 0.60-0.89 & suggest tier & Pass \\
9.4 & Unknown confidence tier & Very short/vague text & tier=unknown, confidence \textless 0.60 & unknown tier & Pass \\
9.5 & Rule-based hackathon override & Complaint mentioning "hackathon" & Category = Hackathon/Event & Override works & Pass \\
9.6 & Negative sentiment detection & Text with angry words & sentiment=Negative, score \textless 0 & Negative detected & Pass \\
9.7 & Positive sentiment detection & Text with appreciative words & sentiment=Positive, score \textgreater 0 & Positive detected & Pass \\
9.8 & Priority boost for negative & Negative sentiment + Medium priority & Priority boosted to High & Boosted & Pass \\
9.9 & Bigram feature extraction & Text with "exam schedule" & bigram "exam\_schedule" in features & Bigram extracted & Pass \\
\end{longtable}

\subsection{Settings and Profile Module}

\begin{longtable}{c p{3.5cm} p{3cm} p{3.5cm} p{3cm} c}
\caption{Settings and Profile Module Test Cases}\\
\toprule
\textbf{TC ID} & \textbf{Description} & \textbf{Input} & \textbf{Expected Output} & \textbf{Actual Output} & \textbf{Status} \\
\midrule
\endfirsthead
\multicolumn{6}{c}{\tablename\ \thetable\ -- Continued} \\
\toprule
\textbf{TC ID} & \textbf{Description} & \textbf{Input} & \textbf{Expected Output} & \textbf{Actual Output} & \textbf{Status} \\
\midrule
\endhead
\bottomrule
\endfoot
10.1 & Update student profile & Change name and phone & Profile updated in DB & Updated & Pass \\
10.2 & Update admin profile & Change name and phone & Profile updated in DB & Updated & Pass \\
10.3 & Change password successfully & Correct current, new password & Password updated & Updated & Pass \\
10.4 & Wrong current password & Incorrect current password & Error: Current password incorrect & Error shown & Pass \\
10.5 & New passwords don't match & Different new and confirm & Error: Passwords do not match & Error shown & Pass \\
10.6 & New password too short & 3-character new password & Error: Minimum 6 characters & Error shown & Pass \\
\end{longtable}

\subsection{CSV Export Module}

\begin{longtable}{c p{3.5cm} p{3cm} p{3.5cm} p{3cm} c}
\caption{CSV Export Module Test Cases}\\
\toprule
\textbf{TC ID} & \textbf{Description} & \textbf{Input} & \textbf{Expected Output} & \textbf{Actual Output} & \textbf{Status} \\
\midrule
\endfirsthead
\multicolumn{6}{c}{\tablename\ \thetable\ -- Continued} \\
\toprule
\textbf{TC ID} & \textbf{Description} & \textbf{Input} & \textbf{Expected Output} & \textbf{Actual Output} & \textbf{Status} \\
\midrule
\endhead
\bottomrule
\endfoot
11.1 & Export all complaints & Click Export CSV button & CSV file downloaded with all complaints & CSV downloaded & Pass \\
11.2 & CSV file format & Check downloaded file & Headers match complaint fields & Headers correct & Pass \\
\end{longtable}

% ============================================================
% CHAPTER 8: CONCLUSION
% ============================================================
\chapter{Conclusion}

\section{Summary}
Complainify is a web-based AI-powered complaint management system that successfully achieves its primary objectives: secure user authentication, complaint submission with automatic AI-based classification, sentiment analysis, priority assignment, role-based dashboards, real-time tracking, email notifications, and CSV report export. The system is developed using the Waterfall model as an SDLC approach.

The Naive Bayes classifier, implemented entirely from scratch, achieves approximately 96\% accuracy on a held-out test dataset of 1,806 samples. The classifier uses a custom stemmer with 25 rules, bigram features, Laplace smoothing, and a three-tier confidence system (Auto, Suggest, Unknown) to indicate classification reliability. The sentiment analyzer, also implemented from scratch using a lexicon-based approach, effectively detects positive, neutral, and negative tones in complaint text and automatically boosts priority for negative-sentiment complaints.

The system provides a modern, responsive user interface built with HTML, CSS, JavaScript, and Tailwind CSS, with separate dashboards for students and administrators. Students can submit complaints, track their status, and manage their profiles. Administrators can manage all complaints, assign them to departments, update statuses, validate complaints, view analytics dashboards with real-time data, generate CSV reports, and send email notifications.

During testing and debugging, all modules were thoroughly tested with comprehensive test cases covering login, registration, forgot password, complaint submission, tracking, admin management, dashboard analytics, email notifications, classifier accuracy, sentiment analysis, profile management, and CSV export. All test cases passed successfully.

\section{Future Enhancements}
Although the system achieves its primary objectives, there are several areas for future improvement:

\begin{itemize}[nolistsep]
    \item \textbf{Multi-language Support:} Adding support for multiple languages to accommodate diverse user populations.
    \item \textbf{Mobile Application:} Developing native mobile applications for iOS and Android platforms.
    \item \textbf{Real-time Push Notifications:} Integrating push notification services for immediate alerts.
    \item \textbf{Advanced AI Models:} Exploring deep learning approaches such as LSTM or BERT for improved classification accuracy.
    \item \textbf{Predictive Analytics:} Implementing predictive models to forecast complaint trends and proactively address recurring issues.
    \item \textbf{Integration with University Systems:} Connecting with existing university management systems for seamless data exchange.
    \item \textbf{Anonymous Complaint Verification:} Implementing mechanisms to verify anonymous complaints while maintaining privacy.
    \item \textbf{Sentiment Model Improvement:} Enhancing the sentiment lexicon with domain-specific terms and contextual understanding.
\end{itemize}

\section{Challenges Faced}
During development, we faced several challenges:

\begin{itemize}[nolistsep]
    \item Implementing the Naive Bayes classifier from scratch required a deep understanding of probability theory, log-space computations to avoid underflow, and Laplace smoothing techniques.
    \item Building a custom stemmer with 25 rules required careful analysis of word patterns in the complaint dataset.
    \item Integrating the sentiment analyzer with the Flask application required careful handling of the import path structure.
    \item The MySQL database schema had to be updated multiple times as new features (sentiment, OTP, assignment tracking) were added.
    \item Debugging AJAX preview functionality for the complaint submission form required careful JavaScript development.
\end{itemize}

Overall, we enjoyed working on this project and gained practical knowledge of implementing core web technologies, machine learning algorithms from scratch, and building a complete full-stack web application. Beyond classroom learning, the project provided us an opportunity to explore and work with various technologies and libraries used during development.

% ============================================================
% REFERENCES
% ============================================================
\begin{thebibliography}{99}
\addcontentsline{toc}{chapter}{References}

\bibitem{deo2025}
Deo, A., Gawande, T., Gore, V., Gawande, R., \& Narode, P. (2025). Smart complaint management system with AI-powered prioritization and escalation. \textit{International Journal of Creative Research Thoughts (IJCRT)}, 13(10), f295--f300.

\bibitem{javed2025}
Javed, R., Wasi, S., Mughal, M. H., \& Bhutto, Z. A. (2025). Machine learning-based smart students' complaint resolution system. \textit{International Journal of Artificial Intelligence \& Mathematical Sciences}, 4(1), 36--55.

\bibitem{prakash2024}
Prakash, J., Swathiramya, R., Balambigai, G., Menaha, R., \& Abhirami, J. S. (2024). AI-driven real-time feedback system for enhanced student support: Leveraging sentiment analysis and machine learning algorithms. \textit{International Journal of Computational and Experimental Science and Engineering}, 10(4).

\bibitem{flask}
Flask Documentation. (2026). Flask Web Framework. Retrieved from \url{https://flask.palletsprojects.com/}

\bibitem{pymysql}
PyMySQL Documentation. (2026). PyMySQL — MySQL driver written in Python. Retrieved from \url{https://pymysql.readthedocs.io/}

\bibitem{pandas}
Pandas Documentation. (2026). pandas: Python Data Analysis Library. Retrieved from \url{https://pandas.pydata.org/}

\bibitem{tailwind}
Tailwind CSS Documentation. (2026). Tailwind CSS: Utility-first CSS framework. Retrieved from \url{https://tailwindcss.com/}

\end{thebibliography}

\end{document}
